\section{强化学习的成就与门槛}

\subsection{Zero 系列的哲学与方法}
Hubert 把 AlphaZero 的成功归结为两个核心：\textbf{search-integrated learning} 与 \textbf{reliable reward from simulation}。在游戏中，模拟环境提供了 perfect feedback，RL 不再依赖人类 label，而是靠 search 的 wins/losses 作为 ground truth。
\begin{knowledgebox}{Zero 哲学的三个支柱}
1）Self-play generating curriculum；2）Search results distilled into policy/value；3）Preference for small, reusable modules rather than monolithic black-box policies。
\end{knowledgebox}

他说：“Zero 系列的辉煌不是单个模型的成功，而是一种微调过的工程流程：每一次 self-play 都会产出 curriculum，policy/value 会被重构成 search-ready 的 heuristic module。”这种哲学直接影响 AlphaProof 的 policy distillation 期望。

\subsection{搜索-学习飞轮}
AlphaZero/AlphaTensor 构建了一个”探索→学习→优化“的飞轮。Hubert 特别指出：如果没有 strong search，就得不到 high-quality trajectories；如果没有 fast learning，就无法 internalize 新知识；如果没有 policy/value 的紧密结合，search 会陷入 thousand-node dead-end。
\begin{warningbox}{飞轮缺口与潜在失败}
1）Search 质量低：agent 无法发现强解，policy 更新 weak；2）带噪声 reward：学习阶段引入 hallucinated trajectories；3）Compute 受限：search depth 不够导致 learning 无 convergence。
\end{warningbox}

Hubert 把这个飞轮比作“在黑洞旁边开挖”：search 不能停留在浅层，否则 learning 永远没得 push；learning 也要及时 absorb search 发现的新 lemma，否则 search 继续重复旧路。AlphaProof 团队用 two-tier policy buffer 让 search-area 与 learning-rate 跟踪，避免出现 drift。

\subsection{探索、反馈与泛化}
AlphaZero 的 feedback 是 deterministic 的 win-rate；AlphaProof 必须面对“formal proof correctness”——即只要 Lean 证明不通过，feedback 就是 fail。Hubert 讲到：“在数学里，reward 要么+1，要么0，no gray area，这迫使 RL 发展 trace-aware regret minimization。” 搜索 vs heuristics 的对比：
\begin{table}[H]
\centering
\begin{tabular}{p{4cm}p{4cm}p{4cm}}
\toprule
维度 & Search-driven RL & Heuristic/Scripted \\
\midrule
反馈质量 & 完美验证，replay buffer 精度高 & 信号噪声高，容易 drift \\
泛化能力 & 通过 policy distillation 实现跨定理迁移 & 需要手工 engineering \\
可解释性 & 可以 replay trajectory，提供 proof term & 难以化简到可检查步骤 \\
\bottomrule
\end{tabular}
\caption{AlphaProof 的 RL 反馈与启发式方法对比}
\end{table}

\subsection{本章小结}
AlphaZero 成功的三个层面（search、feedback、learning）成为 AlphaProof 的起点；必须关注飞轮中的任何断裂、反馈信号的崩盘以及泛化的限制，才能确保 RL agent 在数学世界里稳定进步。
\newpage

